% FILE: CSE-16_TermFinal.tex
\documentclass{article}
\usepackage{amsmath}
\usepackage{amssymb}
\begin{document}

%%%%%%%%%%%%%%%%%%%%%%%%%%%%%%%%%%%%%%%%%%%%%%%%%%
% QUESTION_START
%%%%%%%%%%%%%%%%%%%%%%%%%%%%%%%%%%%%%%%%%%%%%%%%%%
% id=cse16-final-q5a
% discipline=CSE
% year=2016
% exam_type=Term Final
% section=B
% ct_number=UNKNOWN
% question_number=5a
% topic=Limits
% subtopic=General
% difficulty=Easy
% length=Short
% question_type=New
% frequency=1
% appearances=
% OCR_STATUS=VERIFIED
\section*{Term Final (Sec B) - Q5(a)}
Question:
Define integrand, integral and integration. Why does a constant need to be added after integration?

Solution:
Step 1: Define Integrand.
The integrand is the function $f(x)$ that is being integrated. In the mathematical notation $\int f(x) \, dx$, the function $f(x)$ is the integrand.

Step 2: Define Integral.
An integral is the mathematical object resulting from the process of integration. It can be a definite integral (representing a signed area) or an indefinite integral (representing a family of antiderivatives).

Step 3: Define Integration.
Integration is the mathematical operation of calculating an integral. It acts as the inverse process of differentiation.

Step 4: Explain why a constant of integration $C$ must be added.
Let $F(x)$ be an antiderivative of $f(x)$, such that $F'(x) = f(x)$. For any arbitrary constant $C$, the derivative of a constant is zero:
\[
\frac{d}{dx} [F(x) + C] = F'(x) + 0 = f(x)
\]
Because differentiation removes constant terms, the reverse process (indefinite integration) cannot uniquely determine the original constant without additional initial conditions. Thus, to represent the most general family of antiderivatives, we add an arbitrary constant $C$.

Final Answer:
\[
\boxed{\text{Constant } C \text{ represents the lost constant term during differentiation.}}
\]
%%%%%%%%%%%%%%%%%%%%%%%%%%%%%%%%%%%%%%%%%%%%%%%%%%
% QUESTION_END
%%%%%%%%%%%%%%%%%%%%%%%%%%%%%%%%%%%%%%%%%%%%%%%%%%

%%%%%%%%%%%%%%%%%%%%%%%%%%%%%%%%%%%%%%%%%%%%%%%%%%
% QUESTION_START
%%%%%%%%%%%%%%%%%%%%%%%%%%%%%%%%%%%%%%%%%%%%%%%%%%
% id=cse16-final-q5b
% discipline=CSE
% year=2016
% exam_type=Term Final
% section=B
% ct_number=UNKNOWN
% question_number=5b
% topic=Definite Integration
% subtopic=General
% difficulty=Medium
% length=Medium
% question_type=New
% frequency=1
% appearances=
% OCR_STATUS=VERIFIED
\section*{Term Final (Sec B) - Q5(b)}
Question:
Describe the rectangle method for finding area. By using this method find the area of the curve $y = x^2$ over the interval $[0, 1]$.

Solution:
Step 1: Describe the rectangle method (Riemann Sums).
To estimate the area under a continuous curve $y = f(x)$ over an interval $[a, b]$:
1. Divide $[a, b]$ into $n$ equal subintervals of width $\Delta x = \frac{b-a}{n}$.
2. The division points are $x_i = a + i\Delta x$ for $i = 0, 1, \dots, n$.
3. Approximate the area of each subinterval as a rectangle of width $\Delta x$ and height $f(x_i)$.
4. The total area is estimated by the Riemann sum:
   \[
   A \approx \sum_{i=1}^n f(x_i) \Delta x
   \]
5. The exact area is obtained by taking the limit as $n \to \infty$:
   \[
   A = \lim_{n \to \infty} \sum_{i=1}^n f(x_i) \Delta x
   \]

Step 2: Set up the sum for $y = x^2$ over $[0, 1]$.
Here, $f(x) = x^2$, $a = 0$, and $b = 1$.
\[
\Delta x = \frac{1-0}{n} = \frac{1}{n}
\]
Using right endpoints, the evaluation points are $x_i = 0 + i \Delta x = \frac{i}{n}$.
Substituting these into the area limit formula:
\[
A = \lim_{n \to \infty} \sum_{i=1}^n \left(\frac{i}{n}\right)^2 \left(\frac{1}{n}\right) = \lim_{n \to \infty} \frac{1}{n^3} \sum_{i=1}^n i^2
\]

Step 3: Evaluate the sum using the sum of squares formula.
Recall that $\sum_{i=1}^n i^2 = \frac{n(n+1)(2n+1)}{6}$.
\[
A = \lim_{n \to \infty} \frac{n(n+1)(2n+1)}{6n^3}
\]
\[
A = \lim_{n \to \infty} \frac{2n^3 + 3n^2 + n}{6n^3} = \lim_{n \to \infty} \left( \frac{2}{6} + \frac{3}{6n} + \frac{1}{6n^2} \right) = \frac{1}{3}
\]

Final Answer:
\[
\boxed{\frac{1}{3}}
\]
%%%%%%%%%%%%%%%%%%%%%%%%%%%%%%%%%%%%%%%%%%%%%%%%%%
% QUESTION_END
%%%%%%%%%%%%%%%%%%%%%%%%%%%%%%%%%%%%%%%%%%%%%%%%%%

%%%%%%%%%%%%%%%%%%%%%%%%%%%%%%%%%%%%%%%%%%%%%%%%%%
% QUESTION_START
%%%%%%%%%%%%%%%%%%%%%%%%%%%%%%%%%%%%%%%%%%%%%%%%%%
% id=cse16-final-q5c
% discipline=CSE
% year=2016
% exam_type=Term Final
% section=B
% ct_number=UNKNOWN
% question_number=5c
% topic=Definite Integration
% subtopic=General
% difficulty=Easy
% length=Short
% question_type=New
% frequency=1
% appearances=
% OCR_STATUS=VERIFIED
\section*{Term Final (Sec B) - Q5(c)}
Question:
What do you mean by antiderivative and antiderivative method?

Solution:
Step 1: Define Antiderivative.
A function $F(x)$ is called an antiderivative of a function $f(x)$ on an interval $I$ if the derivative of $F(x)$ is equal to $f(x)$ for all $x \in I$:
\[
F'(x) = f(x)
\]

Step 2: Define the Antiderivative Method.
The antiderivative method (also known as the Fundamental Theorem of Calculus, Part 2) is an analytical technique to evaluate definite integrals. Instead of computing Riemann sums directly, we find any antiderivative $F(x)$ of $f(x)$ and evaluate:
\[
\int_a^b f(x) \, dx = F(b) - F(a)
\]

Final Answer:
\[
\boxed{\int_a^b f(x) \, dx = F(b) - F(a)}
\]
%%%%%%%%%%%%%%%%%%%%%%%%%%%%%%%%%%%%%%%%%%%%%%%%%%
% QUESTION_END
%%%%%%%%%%%%%%%%%%%%%%%%%%%%%%%%%%%%%%%%%%%%%%%%%%

%%%%%%%%%%%%%%%%%%%%%%%%%%%%%%%%%%%%%%%%%%%%%%%%%%
% QUESTION_START
%%%%%%%%%%%%%%%%%%%%%%%%%%%%%%%%%%%%%%%%%%%%%%%%%%
% id=cse16-final-q6a
% discipline=CSE
% year=2016
% exam_type=Term Final
% section=B
% ct_number=UNKNOWN
% question_number=6a
% topic=Definite Integration
% subtopic=General
% difficulty=Medium
% length=Medium
% question_type=New
% frequency=1
% appearances=
% OCR_STATUS=VERIFIED
\section*{Term Final (Sec B) - Q6(a)}
Question:
If $f(x)$ is integrable in $(a, b)$, $a < b$ and if there exists a function $F'(x) = f(x)$ in $(a, b)$, then show that $\int_a^b f(x) \, dx = F(b) - F(a)$.

Solution:
Step 1: Set up a partition of $[a, b]$.
Let $P = \{a = x_0, x_1, x_2, \dots, x_n = b\}$ be any partition of the interval $[a, b]$ with subinterval widths $\Delta x_i = x_i - x_{i-1}$.

Step 2: Apply the Mean Value Theorem to $F(x)$.
Since $F(x)$ is differentiable on $(a, b)$ and continuous on $[a, b]$, the Mean Value Theorem applies to $F(x)$ on each subinterval $[x_{i-1}, x_i]$. Thus, there exists some $c_i \in (x_{i-1}, x_i)$ such that:
\[
F(x_i) - F(x_{i-1}) = F'(c_i)(x_i - x_{i-1})
\]
Since $F'(x) = f(x)$:
\[
F(x_i) - F(x_{i-1}) = f(c_i) \Delta x_i
\]

Step 3: Sum over all subintervals.
Summing this relation from $i = 1$ to $n$:
\[
\sum_{i=1}^n [F(x_i) - F(x_{i-1})] = \sum_{i=1}^n f(c_i) \Delta x_i
\]
The left side is a telescoping sum:
\[
[F(x_1) - F(x_0)] + [F(x_2) - F(x_1)] + \dots + [F(x_n) - F(x_{n-1})] = F(x_n) - F(x_0) = F(b) - F(a)
\]
Therefore, we have:
\[
F(b) - F(a) = \sum_{i=1}^n f(c_i) \Delta x_i
\]

Step 4: Take the limit as the norm of partition approaches 0.
Since $f(x)$ is integrable, the Riemann sum on the right converges to the definite integral as $n \to \infty$:
\[
F(b) - F(a) = \lim_{\|P\| \to 0} \sum_{i=1}^n f(c_i) \Delta x_i = \int_a^b f(x) \, dx
\]
Thus, we obtain:
\[
\int_a^b f(x) \, dx = F(b) - F(a)
\]
Hence proved.

Final Answer:
\[
\boxed{\int_a^b f(x) \, dx = F(b) - F(a)}
\]
%%%%%%%%%%%%%%%%%%%%%%%%%%%%%%%%%%%%%%%%%%%%%%%%%%
% QUESTION_END
%%%%%%%%%%%%%%%%%%%%%%%%%%%%%%%%%%%%%%%%%%%%%%%%%%

%%%%%%%%%%%%%%%%%%%%%%%%%%%%%%%%%%%%%%%%%%%%%%%%%%
% QUESTION_START
%%%%%%%%%%%%%%%%%%%%%%%%%%%%%%%%%%%%%%%%%%%%%%%%%%
% id=cse16-final-q6b
% discipline=CSE
% year=2016
% exam_type=Term Final
% section=B
% ct_number=UNKNOWN
% question_number=6b
% topic=Definite Integration
% subtopic=General
% difficulty=Hard
% length=Long
% question_type=New
% frequency=1
% appearances=
% OCR_STATUS=VERIFIED
\section*{Term Final (Sec B) - Q6(b)}
Question:
Find the limit when $n \to \infty$ of the product $\left(1 + \frac{1}{n}\right)\left(1 + \frac{2}{n}\right)^{\frac{1}{2}}\left(1 + \frac{3}{n}\right)^{\frac{1}{3}} \dots \left(1 + \frac{n}{n}\right)^{\frac{1}{n}}$.

Solution:
Step 1: Set up the limit of the product.
Let the limit of the product be $P$:
\[
P = \lim_{n \to \infty} \prod_{r=1}^n \left(1 + \frac{r}{n}\right)^{\frac{1}{r}}
\]
Take the natural logarithm of both sides:
\[
\ln P = \lim_{n \to \infty} \ln \left[ \prod_{r=1}^n \left(1 + \frac{r}{n}\right)^{\frac{1}{r}} \right] = \lim_{n \to \infty} \sum_{r=1}^n \frac{1}{r} \ln \left(1 + \frac{r}{n}\right)
\]

Step 2: Convert the sum to a Riemann sum format.
Multiply and divide the summand by $n$:
\[
\ln P = \lim_{n \to \infty} \sum_{r=1}^n \left( \frac{n}{r} \right) \ln \left(1 + \frac{r}{n}\right) \frac{1}{n} = \lim_{n \to \infty} \sum_{r=1}^n \frac{\ln\left(1 + \frac{r}{n}\right)}{\frac{r}{n}} \frac{1}{n}
\]
This represents a Riemann sum for the continuous function $g(x) = \frac{\ln(1+x)}{x}$ over $[0, 1]$.

Step 3: Convert the limit of sum to a definite integral.
\[
\ln P = \int_0^1 \frac{\ln(1+x)}{x} \, dx
\]

Step 4: Evaluate the definite integral using power series expansion.
Recall the Taylor series expansion of $\ln(1+x)$:
\[
\ln(1+x) = \sum_{k=1}^\infty (-1)^{k-1} \frac{x^k}{k} \implies \frac{\ln(1+x)}{x} = \sum_{k=1}^\infty (-1)^{k-1} \frac{x^{k-1}}{k}
\]
Integrating term-by-term:
\[
\int_0^1 \frac{\ln(1+x)}{x} \, dx = \sum_{k=1}^\infty \frac{(-1)^{k-1}}{k} \left[ \frac{x^k}{k} \right]_0^1 = \sum_{k=1}^\infty \frac{(-1)^{k-1}}{k^2}
\]
This alternating series evaluates to:
\[
\sum_{k=1}^\infty \frac{(-1)^{k-1}}{k^2} = \frac{\pi^2}{12}
\]
Thus:
\[
\ln P = \frac{\pi^2}{12} \implies P = e^{\frac{\pi^2}{12}}
\]

Final Answer:
\[
\boxed{e^{\frac{\pi^2}{12}}}
\]
%%%%%%%%%%%%%%%%%%%%%%%%%%%%%%%%%%%%%%%%%%%%%%%%%%
% QUESTION_END
%%%%%%%%%%%%%%%%%%%%%%%%%%%%%%%%%%%%%%%%%%%%%%%%%%

%%%%%%%%%%%%%%%%%%%%%%%%%%%%%%%%%%%%%%%%%%%%%%%%%%
% QUESTION_START
%%%%%%%%%%%%%%%%%%%%%%%%%%%%%%%%%%%%%%%%%%%%%%%%%%
% id=cse16-final-q7a
% discipline=CSE
% year=2016
% exam_type=Term Final
% section=B
% ct_number=UNKNOWN
% question_number=7a
% topic=Definite Integration
% subtopic=General
% difficulty=Hard
% length=Long
% question_type=Repeated PYQ Pattern
% frequency=1
% appearances=
% OCR_STATUS=VERIFIED
\section*{Term Final (Sec B) - Q7(a)}
Question:
Define Beta and Gamma function. Show that $\int_0^{\pi/2} \sin^m \theta \cos^n \theta \, d\theta = \frac{\Gamma\left(\frac{m+1}{2}\right) \Gamma\left(\frac{n+1}{2}\right)}{2 \Gamma\left(\frac{m+n+2}{2}\right)}$.

Solution:
Step 1: Define the Beta and Gamma functions.
- The Gamma function $\Gamma(z)$ is defined as:
  \[
  \Gamma(z) = \int_0^\infty e^{-t} t^{z-1} \, dt \quad (z > 0)
  \]
- The Beta function $B(a, b)$ is defined as:
  \[
  B(a, b) = \int_0^1 x^{a-1} (1-x)^{b-1} \, dx \quad (a, b > 0)
  \]
- The relation between them is:
  \[
  B(a, b) = \frac{\Gamma(a)\Gamma(b)}{\Gamma(a+b)}
  \]

Step 2: Express Beta function in trigonometric form.
Let $x = \sin^2\theta \implies dx = 2\sin\theta\cos\theta \, d\theta$.
The limits change as follows: when $x=0 \implies \theta=0$; when $x=1 \implies \theta=\pi/2$.
Substitute these into the Beta definition:
\[
B(a, b) = \int_0^{\pi/2} (\sin^2\theta)^{a-1} (1-\sin^2\theta)^{b-1} \cdot 2\sin\theta\cos\theta \, d\theta
\]
\[
B(a, b) = 2 \int_0^{\pi/2} \sin^{2a-1}\theta \cos^{2b-1}\theta \, d\theta
\]

Step 3: Relate exponents to $m$ and $n$.
Let $2a-1 = m \implies a = \frac{m+1}{2}$.
Let $2b-1 = n \implies b = \frac{n+1}{2}$.
Substituting these parameter transformations:
\[
B\left(\frac{m+1}{2}, \frac{n+1}{2}\right) = 2 \int_0^{\pi/2} \sin^m\theta \cos^n\theta \, d\theta
\]
\[
\int_0^{\pi/2} \sin^m\theta \cos^n\theta \, d\theta = \frac{1}{2} B\left(\frac{m+1}{2}, \frac{n+1}{2}\right)
\]

Step 4: Substitute the Gamma function identity.
Using $B(a, b) = \frac{\Gamma(a)\Gamma(b)}{\Gamma(a+b)}$:
\[
\int_0^{\pi/2} \sin^m\theta \cos^n\theta \, d\theta = \frac{\Gamma\left(\frac{m+1}{2}\right)\Gamma\left(\frac{n+1}{2}\right)}{2\Gamma\left(\frac{m+1}{2} + \frac{n+1}{2}\right)} = \frac{\Gamma\left(\frac{m+1}{2}\right)\Gamma\left(\frac{n+1}{2}\right)}{2\Gamma\left(\frac{m+n+2}{2}\right)}
\]
Hence proved.

Final Answer:
\[
\boxed{\int_0^{\pi/2} \sin^m\theta \cos^n\theta \, d\theta = \frac{\Gamma\left(\frac{m+1}{2}\right)\Gamma\left(\frac{n+1}{2}\right)}{2\Gamma\left(\frac{m+n+2}{2}\right)}}
\]
%%%%%%%%%%%%%%%%%%%%%%%%%%%%%%%%%%%%%%%%%%%%%%%%%%
% QUESTION_END
%%%%%%%%%%%%%%%%%%%%%%%%%%%%%%%%%%%%%%%%%%%%%%%%%%

%%%%%%%%%%%%%%%%%%%%%%%%%%%%%%%%%%%%%%%%%%%%%%%%%%
% QUESTION_START
%%%%%%%%%%%%%%%%%%%%%%%%%%%%%%%%%%%%%%%%%%%%%%%%%%
% id=cse16-final-q7b
% discipline=CSE
% year=2016
% exam_type=Term Final
% section=B
% ct_number=UNKNOWN
% question_number=7b
% topic=Definite Integration
% subtopic=General
% difficulty=Medium
% length=Medium
% question_type=New
% frequency=1
% appearances=
% OCR_STATUS=VERIFIED
\section*{Term Final (Sec B) - Q7(b)}
Question:
Show that $\int_0^\infty \frac{x \, dx}{(x+1)(1+x^2)} = \frac{\pi}{4}$.

Solution:
Step 1: Perform partial fraction decomposition.
\[
\frac{x}{(x+1)(1+x^2)} = \frac{A}{x+1} + \frac{Bx+C}{1+x^2}
\]
Multiply both sides by $(x+1)(1+x^2)$:
\[
x = A(1+x^2) + (Bx+C)(x+1)
\]
\[
x = (A+B)x^2 + (B+C)x + (A+C)
\]
Equating coefficients:
- $A+B = 0 \implies B = -A$
- $B+C = 1 \implies -A+C = 1$
- $A+C = 0 \implies C = -A$

Substituting $C = -A$ into the second equation:
\[
-A - A = 1 \implies -2A = 1 \implies A = -\frac{1}{2}
\]
This gives $B = \frac{1}{2}$ and $C = \frac{1}{2}$.

Step 2: Express the integral with the partial fractions.
\[
\int \frac{x \, dx}{(x+1)(1+x^2)} = -\frac{1}{2} \int \frac{dx}{x+1} + \frac{1}{2} \int \frac{x+1}{1+x^2} \, dx
\]
\[
= -\frac{1}{2} \int \frac{dx}{x+1} + \frac{1}{4} \int \frac{2x}{1+x^2} \, dx + \frac{1}{2} \int \frac{dx}{1+x^2}
\]
\[
= -\frac{1}{2}\ln|x+1| + \frac{1}{4}\ln(1+x^2) + \frac{1}{2}\tan^{-1} x + C
\]
Combine the logarithmic terms:
\[
= \frac{1}{4}\ln\left( \frac{1+x^2}{(x+1)^2} \right) + \frac{1}{2}\tan^{-1} x + C
\]

Step 3: Evaluate limits from $0$ to $\infty$.
- For the lower limit $x = 0$:
  \[
  \frac{1}{4}\ln(1) + \frac{1}{2}\tan^{-1}(0) = 0
  \]
- For the upper limit as $x \to \infty$:
  \[
  \lim_{x \to \infty} \frac{1+x^2}{(x+1)^2} = \lim_{x \to \infty} \frac{1+x^2}{x^2+2x+1} = 1 \implies \ln(1) = 0
  \]
  \[
  \lim_{x \to \infty} \frac{1}{2}\tan^{-1} x = \frac{1}{2}\left(\frac{\pi}{2}\right) = \frac{\pi}{4}
  \]
Subtracting the evaluations:
\[
\left[ \frac{1}{4}\ln\left(\frac{1+x^2}{(x+1)^2}\right) + \frac{1}{2}\tan^{-1} x \right]_0^\infty = \frac{\pi}{4} - 0 = \frac{\pi}{4}
\]
Hence proved.

Final Answer:
\[
\boxed{\frac{\pi}{4}}
\]
%%%%%%%%%%%%%%%%%%%%%%%%%%%%%%%%%%%%%%%%%%%%%%%%%%
% QUESTION_END
%%%%%%%%%%%%%%%%%%%%%%%%%%%%%%%%%%%%%%%%%%%%%%%%%%

%%%%%%%%%%%%%%%%%%%%%%%%%%%%%%%%%%%%%%%%%%%%%%%%%%
% QUESTION_START
%%%%%%%%%%%%%%%%%%%%%%%%%%%%%%%%%%%%%%%%%%%%%%%%%%
% id=cse16-final-q8ai
% discipline=CSE
% year=2016
% exam_type=Term Final
% section=B
% ct_number=UNKNOWN
% question_number=8ai
% topic=Indefinite Integration
% subtopic=Integration by Parts
% difficulty=Medium
% length=Medium
% question_type=New
% frequency=1
% appearances=
% OCR_STATUS=VERIFIED
\section*{Term Final (Sec B) - Q8(a)(i)}
Question:
Evaluate:
\[
\int e^{ax} \sin bx \, dx
\]

Solution:
Step 1: Set up integration by parts.
Let $I = \int e^{ax} \sin bx \, dx$.
Applying integration by parts with:
- $u = \sin bx \implies du = b\cos bx \, dx$
- $dv = e^{ax} dx \implies v = \frac{1}{a} e^{ax}$

\[
I = \frac{1}{a} e^{ax} \sin bx - \frac{b}{a} \int e^{ax} \cos bx \, dx
\]

Step 2: Apply integration by parts a second time.
For the integral $\int e^{ax} \cos bx \, dx$, let:
- $u_2 = \cos bx \implies du_2 = -b\sin bx \, dx$
- $dv_2 = e^{ax} dx \implies v_2 = \frac{1}{a} e^{ax}$

\[
\int e^{ax} \cos bx \, dx = \frac{1}{a} e^{ax} \cos bx + \frac{b}{a} \int e^{ax} \sin bx \, dx
\]
Substituting this back into the expression for $I$:
\[
I = \frac{1}{a} e^{ax} \sin bx - \frac{b}{a} \left[ \frac{1}{a} e^{ax} \cos bx + \frac{b}{a} I \right]
\]
\[
I = \frac{1}{a} e^{ax} \sin bx - \frac{b}{a^2} e^{ax} \cos bx - \frac{b^2}{a^2} I
\]

Step 3: Solve for $I$.
\[
I \left( 1 + \frac{b^2}{a^2} \right) = \frac{e^{ax}}{a^2} (a\sin bx - b\cos bx)
\]
\[
I \left( \frac{a^2+b^2}{a^2} \right) = \frac{e^{ax}}{a^2} (a\sin bx - b\cos bx)
\]
\[
I = \frac{e^{ax}}{a^2+b^2} (a\sin bx - b\cos bx) + C
\]

Final Answer:
\[
\boxed{\frac{e^{ax}}{a^2+b^2} (a\sin bx - b\cos bx) + C}
\]
%%%%%%%%%%%%%%%%%%%%%%%%%%%%%%%%%%%%%%%%%%%%%%%%%%
% QUESTION_END
%%%%%%%%%%%%%%%%%%%%%%%%%%%%%%%%%%%%%%%%%%%%%%%%%%

%%%%%%%%%%%%%%%%%%%%%%%%%%%%%%%%%%%%%%%%%%%%%%%%%%
% QUESTION_START
%%%%%%%%%%%%%%%%%%%%%%%%%%%%%%%%%%%%%%%%%%%%%%%%%%
% id=cse16-final-q8aii
% discipline=CSE
% year=2016
% exam_type=Term Final
% section=B
% ct_number=UNKNOWN
% question_number=8aii
% topic=Indefinite Integration
% subtopic=Substitution
% difficulty=Medium
% length=Medium
% question_type=New
% frequency=1
% appearances=
% OCR_STATUS=VERIFIED
\section*{Term Final (Sec B) - Q8(a)(ii)}
Question:
Evaluate:
\[
\int \frac{x^{1/2}}{1+x^{3/4}} \, dx
\]

Solution:
Step 1: Set up fractional exponent substitution.
Let $x = u^4 \implies dx = 4u^3 \, du$.
Then $x^{1/2} = (u^4)^{1/2} = u^2$ and $x^{3/4} = (u^4)^{3/4} = u^3$.
Substitute these into the integral:
\[
\int \frac{u^2}{1+u^3} (4u^3 \, du) = 4 \int \frac{u^5}{1+u^3} \, du
\]

Step 2: Perform polynomial long division.
\[
\frac{u^5}{u^3+1} = u^2 - \frac{u^2}{u^3+1}
\]
Thus, the integral becomes:
\[
4 \int \left( u^2 - \frac{u^2}{u^3+1} \right) du = 4 \left[ \frac{u^3}{3} - \frac{1}{3}\ln|u^3+1| \right] + C
\]
\[
= \frac{4}{3} u^3 - \frac{4}{3} \ln|u^3+1| + C
\]

Step 3: Substitute back $u = x^{1/4} \implies u^3 = x^{3/4}$.
\[
= \frac{4}{3} x^{3/4} - \frac{4}{3} \ln(1+x^{3/4}) + C
\]

Final Answer:
\[
\boxed{\frac{4}{3} x^{3/4} - \frac{4}{3} \ln(1+x^{3/4}) + C}
\]
%%%%%%%%%%%%%%%%%%%%%%%%%%%%%%%%%%%%%%%%%%%%%%%%%%
% QUESTION_END
%%%%%%%%%%%%%%%%%%%%%%%%%%%%%%%%%%%%%%%%%%%%%%%%%%

%%%%%%%%%%%%%%%%%%%%%%%%%%%%%%%%%%%%%%%%%%%%%%%%%%
% QUESTION_START
%%%%%%%%%%%%%%%%%%%%%%%%%%%%%%%%%%%%%%%%%%%%%%%%%%
% id=cse16-final-q8aiii
% discipline=CSE
% year=2016
% exam_type=Term Final
% section=B
% ct_number=UNKNOWN
% question_number=8aiiii
% topic=Indefinite Integration
% subtopic=General
% difficulty=Hard
% length=Long
% question_type=New
% frequency=1
% appearances=
% OCR_STATUS=VERIFIED
\section*{Term Final (Sec B) - Q8(a)(iii)}
Question:
Evaluate:
\[
\int \frac{dx}{(x-a)\sqrt{(x-a)(b-x)}}
\]

Solution:
Step 1: Set up fractional substitution.
Let $x-a = \frac{1}{u} \implies dx = -\frac{1}{u^2} \, du$.
Also, rewrite the expression inside the radical:
\[
b-x = b - \left(a + \frac{1}{u}\right) = (b-a) - \frac{1}{u}
\]
Substituting these into the radical term:
\[
\sqrt{(x-a)(b-x)} = \sqrt{\frac{1}{u} \left( (b-a) - \frac{1}{u} \right)} = \frac{1}{u}\sqrt{(b-a)u - 1}
\]

Step 2: Perform the full algebraic substitution.
\[
\int \frac{-\frac{1}{u^2} \, du}{\frac{1}{u} \cdot \frac{1}{u}\sqrt{(b-a)u - 1}} = -\int \frac{du}{\sqrt{(b-a)u - 1}}
\]

Step 3: Integrate the simplified expression.
Let $w = (b-a)u - 1 \implies dw = (b-a) \, du$.
\[
-\frac{1}{b-a} \int w^{-1/2} \, dw = -\frac{2}{b-a} w^{1/2} + C = -\frac{2}{b-a}\sqrt{(b-a)u - 1} + C
\]

Step 4: Substitute back $u = \frac{1}{x-a}$.
\[
= -\frac{2}{b-a}\sqrt{ \frac{b-a}{x-a} - 1 } + C
\]
\[
= -\frac{2}{b-a}\sqrt{ \frac{b-a - x + a}{x-a} } + C = -\frac{2}{b-a}\sqrt{\frac{b-x}{x-a}} + C
\]

Final Answer:
\[
\boxed{-\frac{2}{b-a}\sqrt{\frac{b-x}{x-a}} + C}
\]
%%%%%%%%%%%%%%%%%%%%%%%%%%%%%%%%%%%%%%%%%%%%%%%%%%
% QUESTION_END
%%%%%%%%%%%%%%%%%%%%%%%%%%%%%%%%%%%%%%%%%%%%%%%%%%

%%%%%%%%%%%%%%%%%%%%%%%%%%%%%%%%%%%%%%%%%%%%%%%%%%
% QUESTION_START
%%%%%%%%%%%%%%%%%%%%%%%%%%%%%%%%%%%%%%%%%%%%%%%%%%
% id=cse16-final-q8aiv
% discipline=CSE
% year=2016
% exam_type=Term Final
% section=B
% ct_number=UNKNOWN
% question_number=8aiv
% topic=Definite Integration
% subtopic=General
% difficulty=Medium
% length=Medium
% question_type=New
% frequency=1
% appearances=
% OCR_STATUS=VERIFIED
\section*{Term Final (Sec B) - Q8(a)(iv)}
Question:
Evaluate:
\[
\int_0^{\pi/2} \frac{dx}{3+5\cos x}
\]

Solution:
Step 1: Apply Weierstrass half-angle substitution.
Let $t = \tan\left(\frac{x}{2}\right) \implies dx = \frac{2\,dt}{1+t^2}$ and $\cos x = \frac{1-t^2}{1+t^2}$.
Evaluate the boundaries: when $x=0 \implies t=0$; when $x=\pi/2 \implies t=1$.

Step 2: Rewrite the integral.
\[
\int_0^1 \frac{\frac{2\,dt}{1+t^2}}{3 + 5\left(\frac{1-t^2}{1+t^2}\right)} = \int_0^1 \frac{2\,dt}{3(1+t^2) + 5(1-t^2)}
\]
\[
= \int_0^1 \frac{2\,dt}{8 - 2t^2} = \int_0^1 \frac{dt}{4 - t^2}
\]

Step 3: Integrate using standard log formula.
Recall the standard integration identity:
\[
\int \frac{dt}{a^2-t^2} = \frac{1}{2a}\ln\left|\frac{a+t}{a-t}\right|
\]
Here, $a = 2$:
\[
\left[ \frac{1}{4}\ln\left|\frac{2+t}{2-t}\right| \right]_0^1 = \frac{1}{4}\ln(3) - \frac{1}{4}\ln(1) = \frac{1}{4}\ln(3)
\]

Final Answer:
\[
\boxed{\frac{1}{4}\ln 3}
\]
%%%%%%%%%%%%%%%%%%%%%%%%%%%%%%%%%%%%%%%%%%%%%%%%%%
% QUESTION_END
%%%%%%%%%%%%%%%%%%%%%%%%%%%%%%%%%%%%%%%%%%%%%%%%%%

%%%%%%%%%%%%%%%%%%%%%%%%%%%%%%%%%%%%%%%%%%%%%%%%%%
% QUESTION_START
%%%%%%%%%%%%%%%%%%%%%%%%%%%%%%%%%%%%%%%%%%%%%%%%%%
% id=cse16-final-q8av
% discipline=CSE
% year=2016
% exam_type=Term Final
% section=B
% ct_number=UNKNOWN
% question_number=8av
% topic=Definite Integration
% subtopic=General
% difficulty=Medium
% length=Medium
% question_type=New
% frequency=1
% appearances=
% OCR_STATUS=VERIFIED
\section*{Term Final (Sec B) - Q8(a)(v)}
Question:
Evaluate:
\[
\int_0^{\pi/2} \sin\phi \cos\phi \sqrt{a^2 \sin^2\phi + b^2 \cos^2\phi} \, d\phi
\]

Solution:
Step 1: Set up u-substitution.
Let $u = a^2 \sin^2\phi + b^2 \cos^2\phi$.
Differentiating with respect to $\phi$:
\[
du = \left(2a^2 \sin\phi\cos\phi - 2b^2 \sin\phi\cos\phi\right) d\phi = 2(a^2-b^2)\sin\phi\cos\phi \, d\phi
\]
\[
\sin\phi\cos\phi \, d\phi = \frac{du}{2(a^2-b^2)}
\]

Step 2: Update limits of integration.
- When $\phi = 0 \implies u = b^2$.
- When $\phi = \pi/2 \implies u = a^2$.

Step 3: Evaluate the integration.
\[
\int_{b^2}^{a^2} \sqrt{u} \frac{du}{2(a^2-b^2)} = \frac{1}{2(a^2-b^2)} \left[ \frac{2}{3}u^{3/2} \right]_{b^2}^{a^2}
\]
\[
= \frac{1}{3(a^2-b^2)} \left( (a^2)^{3/2} - (b^2)^{3/2} \right) = \frac{a^3 - b^3}{3(a^2-b^2)}
\]
Factor the numerator and the denominator:
\[
= \frac{(a-b)(a^2+ab+b^2)}{3(a-b)(a+b)} = \frac{a^2+ab+b^2}{3(a+b)}
\]

Final Answer:
\[
\boxed{\frac{a^2+ab+b^2}{3(a+b)}}
\]
%%%%%%%%%%%%%%%%%%%%%%%%%%%%%%%%%%%%%%%%%%%%%%%%%%
% QUESTION_END
%%%%%%%%%%%%%%%%%%%%%%%%%%%%%%%%%%%%%%%%%%%%%%%%%%

\end{document}